\section*{Direct proof for Thompson's group (Proposition 21)}
One can also prove the uniform SENS property for Thompson's group $F$ directly. Recall the infinite presentation $F=\langle x_0,x_1,x_2,\ldots\mid x_k^{x_i}=x_{k+1}\ (i<k)\rangle$.
\begin{authorclaim}{Proposition 21}
Let $g=x_3x_2^{-1}\in F$ and define $c_v$ for $v\in\{0,1\}^*$ inductively via
\[
c_\varepsilon=\varepsilon,\qquad c_{v0}=x_1c_v,\qquad c_{v1}=x_0^{-1}x_1c_v.
\]
Finally, for $d\in\mathbb N$ and $v\in\{0,1\}^{\le d}$ let $g_{d,v}=g^{c_v}$. Then $g_{d,v}=[g_{d,v0},g_{d,v1}]$ for all $d\in\mathbb N$ and $v\in\{0,1\}^{\le d}$ and $g_\varepsilon=g\ne1$ in $F$. In particular, $G$ is uniformly SENS.
\end{authorclaim}
\begin{proof}
Obviously, we have $g_\varepsilon=g$ in $F$. Moreover, since $x_2\ne x_3$ (which follows directly from the normal form theorem for $F$; see e.g.\ [12, Corollary 2.7]), we have $g\ne1$ in $F$ (Condition (c) of Definition 1). The identity $g=[g,g^{x_0^{-1}}]^{x_1}=[g^{x_1},g^{x_0^{-1}x_1}]$ is straightforward to check. Thus, we obtain Condition (b):
\[
g_{d,v}=g^{c_v}=[g^{x_1c_v},g^{x_0^{-1}x_1c_v}]
=[g^{c_{v0}},g^{c_{v1}}]=[g_{d,v0},g_{d,v1}].
\]
Moreover, since $|g_{d,v}|=|g|+2|c_v|\le|g|+4d$, we can pad with the identity symbol writing $g_{v,d}$ as a word of length $2^{\mu d}$ for some proper constant $\mu\in\mathbb N$ in order to meet Condition (a) (be aware that the lengths need to be computed over a finite standard generating set $\Sigma$, e.g.\ $\Sigma=\{1,x_0,x_0^{-1},x_1,x_1^{-1}\}$). This shows that $F$ is SENS.
Condition (d) of Definition 1 is also straightforward to check by introducing some more padding: we slightly change the definition of the $c_v$ by setting $c_\varepsilon=\varepsilon$, $c_{v0}=1x_1c_v$, and $c_{v1}=x_0^{-1}x_1c_v$ where $1\in\Sigma$. This new $c_v$ represents the same group element as the old $c_v$, but now we have $|c_v|=2|v|$ for all $v\in\{0,1\}^*$ and all letters at even positions are $x_1$, while letters at odd positions are either $1$ or $x_0^{-1}$ (the $(2j+1)$-st letter of $c_v$ depends on the $(|v|-j)$-th bit of $v$). Notice that we have $|g_{d,v}|=|g|+2|c_v|=|g|+4d$.
On input of $v\in\{0,1\}^d$, $i\in\mathbb N$ (encoded with $\mu d$ bits), and $a\in\Sigma$, first decide whether $1\le i\le2d$, or $2d<i\le2d+|g|$, or $2d+|g|<i\le4d+|g|$, or $i>4d+|g|$. This clearly can be done in DLINTIME (recall that $g$ is a constant). Now assume that $1\le i\le2d$ (i.e., $i$ points into the leading $c_v^{-1}$ of $g_{d,v}=c_v^{-1}gc_v$). If $i$ is odd, one accepts if and only if $a=x_1^{-1}$; if $i$ is even, one accepts if and only if the $i/2$-th bit of $v$ is zero and $a=1$ or if the $i/2$-th bit of $v$ is one and $a=x_0$. The other cases are similar.
\end{proof}
